\section{Reduction to basic algebras and graded reconstruction}
\label{sec:graded-category}

We begin by reducing the existence of a normed multiplicative basis to the
case of basic algebras. We then pass from the category of indecomposable
modules to its radical grading and reconstruct the standard form $A_0$ from
the graded projective objects. Both constructions, together with the
normalization of $A_0$ in the next two sections, apply in arbitrary
characteristic.

\subsection{Multiplicative bases and matrix units}
\label{subsec:nonbasic-basis}

Using matrix units to account for the multiplicities of indecomposable
projective summands, a normed multiplicative basis of a basic algebra lifts
immediately to any Morita-equivalent finite-dimensional algebra.

\begin{lemma}
\label{lem:nonbasic-basis}
Let $A$ be a finite-dimensional algebra over an arbitrary field, and let
$B_0$ be its basic algebra. If $B_0$ has a normed multiplicative basis, then
so does $A$.
\end{lemma}
\begin{proof}
Let $\{e_i\}$ be the complete set of primitive orthogonal idempotents
contained in the given normed multiplicative basis of $B_0$, and let $n_i$
be the multiplicity of the corresponding indecomposable projective module in
the right regular module $A_A$. Every element $b$ of the basis of $B_0$ lies
in a unique corner $e_i B_0 e_j$: indeed, in the decomposition
$b=\sum_{i,j} e_i b e_j$, each nonzero summand belongs to the basis by
multiplicativity, and since basis elements in distinct corners are linearly
independent, only one summand can be nonzero.

Under Morita equivalence, the regular module $A_A$ corresponds to the
$B_0$-module $\bigoplus_i (e_i B_0)^{n_i}$. Taking the opposite of its
endomorphism algebra under our left-to-right composition convention
identifies $A$ with the block matrix algebra whose $(i,j)$-block is
$\operatorname{Mat}_{n_i\times n_j}(e_i B_0 e_j)$. For each basis element
$b\in e_i B_0 e_j$, let $E_{i\alpha,j\beta}(b)$ ($1\leq\alpha\leq n_i$,
$1\leq\beta\leq n_j$) denote the matrix having $b$ in entry $(\alpha,\beta)$
of the $(i,j)$-block and zero elsewhere. These matrices form a basis of $A$
and multiply according to the rule
\[
 E_{i\alpha,j\beta}(b)E_{l\gamma,m\delta}(c)
 =\begin{cases}
 E_{i\alpha,m\delta}(bc),&\text{if }(j,\beta)=(l,\gamma),\\
 0,&\text{otherwise}.
 \end{cases}
\]
Since the given basis of $B_0$ is multiplicative, each product is either
zero or another basis element. Moreover, the diagonal elements
$E_{i\alpha,i\alpha}(e_i)$ form a complete set of primitive orthogonal
idempotents of $A$, and for each $n\geq1$ the $(i,j)$-block of $(\rad A)^n$
is $\operatorname{Mat}_{n_i\times n_j}(e_i(\rad B_0)^n e_j)$. Thus the
constructed basis intersects each radical power $(\rad A)^n$ in a basis.
\end{proof}

From now on we may therefore assume that $A$ is basic, invoking
Lemma~\ref{lem:nonbasic-basis} at the conclusion of each characteristic
branch to recover the general case.

\subsection{Radical layers and mesh relations}

We filter the morphisms in $\ind(\modu A)$ by powers of the categorical
radical, that is, by sums of composites of nonisomorphisms between
indecomposable modules. Because intermediate modules in such a composite need
not be projective, the induced filtration on morphisms between indecomposable
projective modules is in general finer than the ordinary radical filtration
of $A$.

Let $A$ be a basic representation-finite algebra over $k$, and let $\JJ$
denote the categorical Jacobson radical of $\ind(\modu A)$. Since $A$ is
representation-finite, the direct sum of a representative set of
indecomposable modules is a finite-dimensional additive generator; its
endomorphism algebra is finite-dimensional, and $\JJ$ consists of the matrix
entries of its Jacobson radical, so $\JJ$ is nilpotent. We define the
\emph{radical grading} of $\ind(\modu A)$ to be the graded category $\C$
with the same objects as $\ind(\modu A)$ and morphism spaces
\[
 \C(X,Y)=\bigoplus_{n\geq0}\JJ^n(X,Y)/\JJ^{n+1}(X,Y),
\]
writing $\C_n(X,Y)=\JJ^n(X,Y)/\JJ^{n+1}(X,Y)$ for the homogeneous component of
degree $n$. For a morphism $f\in \JJ^n(X,Y)\setminus \JJ^{n+1}(X,Y)$, its
\emph{initial form} is its residue class in $\C_n(X,Y)$. Composition in $\C$
is induced by composition in $\ind(\modu A)$. Since the endomorphism ring of
each indecomposable module is local with residue field $k$, we have
$\C_0(X,X)=k\cdot 1_X$ and $\C_0(X,Y)=0$ for $X\ne Y$.

Let us recall how the Auslander--Reiten quiver $\Gamma_A$ and its mesh
relations arise in this setting. The vertices of $\Gamma_A$ are the objects
of $\ind(\modu A)$, and the arrows $X\to Y$ index a chosen basis of the
space $\Irr(X,Y)=\JJ(X,Y)/\JJ^2(X,Y)$ of irreducible morphisms. For each
nonprojective indecomposable module $Y$, choose an almost split sequence
\[
 0\longrightarrow\tau Y\longrightarrow E_Y
   \longrightarrow Y\longrightarrow0.
\]
Each indecomposable direct summand $X$ of $E_Y$ contributes a component map
$X\to Y$ and a component map $\tau Y\to X$; modulo $\JJ^2$, the former
components form a basis of $\Irr(X,Y)$ and the latter a basis of
$\Irr(\tau Y,X)$ indexed by the multiplicity of $X$ in $E_Y$. The next
well-known lemma ensures that each multiplicity is at most one, so that these
components differ from the chosen arrows only by nonzero scalars.

\begin{lemma}
\label{lem:square-free}
Let $A$ be a basic representation-finite algebra over an algebraically
closed field. Its Auslander--Reiten quiver $\Gamma_A$ has no loops and at
most one arrow between each ordered pair of vertices.
\end{lemma}
\begin{proof}
Consider any irreducible morphism $f:X\to Y$ between indecomposable modules
and factor $f$ through its image. By irreducibility, either the canonical
surjection $X\twoheadrightarrow\im f$ is a split monomorphism or the
inclusion $\im f\hookrightarrow Y$ is a split epimorphism; thus one of the
two factors is an isomorphism, and $f$ is either a proper monomorphism or a
proper epimorphism. In particular, $\dim X\ne\dim Y$, so $\Gamma_A$ has no
loops.

Suppose for a contradiction that there are at least two arrows $X\to Y$ in
$\Gamma_A$. If $\dim X>\dim Y$, any irreducible morphism $X\to Y$ is an
epimorphism, so $Y$ cannot be projective (otherwise the epimorphism would
split). The middle term of the almost split sequence ending at $Y$ then
contains $X^{\oplus 2}$ as a direct summand, which gives
\[
 \dim\tau Y\geq2\dim X-\dim Y>\dim X,
\]
and the first map of the sequence provides at least two arrows
$\tau Y\to X$. Iterating this step produces indecomposable modules of
strictly increasing dimensions, contradicting the assumption that $A$ is
representation-finite. If $\dim X<\dim Y$, the dual argument applied to the
almost split sequence starting at $X$ yields the same contradiction.
\end{proof}

Thanks to Lemma~\ref{lem:square-free}, we can formulate the mesh category of
$\Gamma_A$ using scalar coefficients rather than matrices.

\begin{definition}
\label{def:translation-quiver}
A \emph{translation quiver} is a quiver $Q$ equipped with a bijection $\tau$
from a subset of vertices, called \emph{nonprojective}, onto a subset of
vertices, called \emph{noninjective}, together with a bijection $\sigma$
from the set of arrows ending at nonprojective vertices onto the set of
arrows starting at noninjective vertices, such that $\sigma a:\tau x\to y$
for every arrow $a:y\to x$ with $x$ nonprojective. We call $\sigma a$ the
\emph{mate} of $a$. Vertices outside the domain (respectively, image) of
$\tau$ are called \emph{projective} (respectively, \emph{injective}). When
both $x$ and $y$ are nonprojective, we write
$\tau a=\sigma^2a:\tau y\to\tau x$. The translation quiver $Q$ is
\emph{stable} if every vertex is both nonprojective and noninjective.
\end{definition}

\begin{definition}
\label{def:scalar-mesh-category}
Let $Q$ be a translation quiver with at most one arrow between each ordered
pair of vertices. Given a family $c=(c_{x,a})$ of nonzero scalars
$c_{x,a}\in k^*$ indexed by nonprojective vertices $x$ and arrows $a:y\to x$
ending at $x$, the associated \emph{scalar mesh category} is the path-length
graded quotient category
\[
 C(Q,c)=kQ/\langle r_x:x\text{ nonprojective}\rangle,
 \qquad r_x=\sum_{a:y\to x}c_{x,a}(\sigma a)a,
\]
where $kQ$ denotes the $k$-linear path category of $Q$.
\end{definition}

Multiplying a mesh relation $r_x$ by a nonzero scalar leaves the relation
ideal unchanged, while rescaling the arrows of $Q$ rescales the coefficients
$c_{x,a}$ and induces a graded isomorphism between the corresponding scalar
mesh categories.

For $Q=\Gamma_A$, choose a basis element in each one-dimensional space
$\Irr(X,Y)$ to serve as the unique arrow $a:X\to Y$. For each nonprojective
indecomposable module $Y$, the components of the first and second maps of
the chosen almost split sequence ending at $Y$ take the form
$u_a(\sigma a)$ and $v_a a$ modulo $\JJ^2$, with $u_a,v_a\in k^*$. Since the
composite of the two maps in the almost split sequence is zero, passing to
degree two in $\C$ yields the mesh relation
\begin{equation}
 r_Y=\sum_{a:X\to Y}c_{Y,a}(\sigma a)a,
 \qquad c_{Y,a}=u_av_a.
 \label{eq:scalar-mesh}
\end{equation}
We write $k(\Gamma_A)$ for the scalar mesh category $C(\Gamma_A,c)$ defined
by these coefficients.

Up to isomorphism fixing the objects, $k(\Gamma_A)$ is independent of both
the chosen arrow representatives and the chosen almost split sequences.
Indeed, any isomorphism between two almost split sequences ending at $Y$ acts
by scalars $\xi,\eta\in k^*$ on the indecomposable end terms $\tau Y$ and
$Y$. Because the indecomposable summands of the middle term are pairwise
nonisomorphic by Lemma~\ref{lem:square-free}, the automorphism of the middle
term acts by scalars $\lambda_a\in k^*$ on those summands modulo radical
maps, whose composites with the maps of the sequence vanish modulo $\JJ^2$.
The coefficients therefore transform according to
\[
 u'_a=\lambda_a u_a\xi^{-1},\qquad
 v'_a=\eta v_a\lambda_a^{-1},\qquad
 c'_{Y,a}=\eta\xi^{-1}c_{Y,a},
\]
which simply rescales the entire relation $r_Y$ by $\eta\xi^{-1}$. Moreover,
Lemma~\ref{lem:unit-mesh-coefficients} will show that the arrows can always
be rescaled so that $c_{Y,a}=1$ for all $Y$ and $a$.

In Sections~\ref{sec:scalar-normalization} and~\ref{sec:mesh-deformations}
we will restrict a scalar mesh category $C(Q,c)$ to a subset $S$ of vertices
that is closed under $\tau$ and $\tau^{-1}$ wherever defined (such as the
set of $\tau$-periodic vertices). Let $Q_S$ be the full translation
subquiver on $S$. Quotienting $C(Q,c)$ by the ideal generated by the
identity morphisms $1_x$ for $x\notin S$ and restricting to the objects in
$S$ gives a canonical isomorphism
\begin{equation}
 \left(C(Q,c)/\langle 1_x:x\notin S\rangle\right)\big|_S
       \cong C(Q_S,c|_S).
 \label{eq:mesh-object-quotient}
\end{equation}
Indeed, in the path category $kQ$, killing $1_x$ for $x\notin S$ annihilates
precisely those paths that pass through a vertex outside $S$. Under this
quotient, a mesh relation $r_x$ with $x\notin S$ maps to zero, while a mesh
relation $r_x$ with $x\in S$ retains precisely the terms $(\sigma a)a$ whose
middle vertex lies in $S$, which is the mesh relation of $Q_S$ at $x$. In
particular, if $C(Q,c)$ is finite-dimensional, then so is $C(Q_S,c|_S)$.

\subsection{Exactness on radical layers}

We next show that almost split sequences remain exact after passing to the
associated graded category $\C$. This exactness is the key to identifying
$\C$ with the mesh category $k(\Gamma_A)$ in
Proposition~\ref{prop:graded-realization}(3) and to reconstructing the module
category $\modu A_0$ in Proposition~\ref{prop:projective-evaluation}. Our
argument adapts Iyama's complementary-monomorphism technique for
$\tau$-categories \cite[Sections~3.3, 3.6 and~4.2]{IyamaI} directly to module
categories.

\begin{lemma}
\label{lem:strict-radical-layers}
Let $A$ be an Artin algebra, let $\JJ$ be the categorical radical of its
category of finite-length modules, and let $a:U\to Y$ be a radical
monomorphism. Suppose that $a+h$ is a monomorphism for every
$h\in \JJ^2(U,Y)$. Then for every finite-length module $W$ and every
$n\geq1$,
\[
             a\circ z\in \JJ^n(W,Y)\quad\Longleftrightarrow\quad
             z\in \JJ^{n-1}(W,U).
\]
Furthermore, any finite direct sum of left almost split monomorphisms
satisfies the hypothesis that all such $\JJ^2$-perturbations remain
monomorphisms.
\end{lemma}
\begin{proof}
\emph{Step 1: Almost split kernels and the perturbation property.}
For any finite direct sum $Y$ of indecomposable modules, let $p:E_Y\to Y$ be
the direct sum of the minimal right almost split morphisms into its
indecomposable summands (using the radical inclusion $\rad P\hookrightarrow P$
at projective summands). By the lifting property of right almost split maps,
every radical morphism into $Y$ factors through $p$. For $n\geq0$, writing a
morphism in $\JJ^{n+1}(W,Y)$ as a sum of composites of $n+1$ radical morphisms
and factoring the last radical morphism through $p$ shows that
\begin{equation}
 p\circ \JJ^n(W,E_Y)=\JJ^{n+1}(W,Y)\qquad(n\geq0).
 \label{eq:right-radical-surjection}
\end{equation}
At each nonprojective indecomposable summand of $Y$, the kernel of $p$ is
the left almost split morphism of the corresponding almost split sequence,
while at each projective summand the inclusion $\rad P\hookrightarrow P$ has
zero kernel. Consequently, the kernel inclusion $i:T\to E_Y$ of $p$ is a
finite direct sum of left almost split monomorphisms.

More generally, let $i=\bigoplus_j i_j:T\to E$ be any finite direct sum of
left almost split monomorphisms $i_j:T_j\to E_j$. Any morphism in
$\JJ^2(T,Z)$ is a sum of composites $v\circ u$ with $u$ and $v$ radical. Since
the restriction of $u$ to each $T_j$ is radical, it factors through $i_j$ as
$u_j\circ i_j$, and $v\circ u_j$ remains radical. Thus
\[
 \JJ^2(T,Z)=\{q\circ i:q\in \JJ(E,Z)\}.
\]
Taking $Z=E$, any $h\in \JJ^2(T,E)$ has the form $h=q\circ i$ with
$q\in \JJ(E,E)$. Since $1+q$ is an automorphism of $E$, the perturbation
$i+h=(1+q)\circ i$ is again a monomorphism.

\emph{Step 2: Splitting off the given monomorphism.}
Returning to the given monomorphism $a:U\to Y$, since $a$ is radical and
$p:E_Y\to Y$ is right almost split, we can write $a=p\circ u$ for some
$u:U\to E_Y$. By performing row and column operations to split off any
invertible components of $u$ between indecomposable summands, we obtain
decompositions $U=U''\oplus U'$ and $E_Y=U''\oplus E'$ under which
$u=1_{U''}\oplus u'$ with $u'\in \JJ(U',E')$. Since $p$ is radical, the
restriction $a|_{U'}=p|_{E'}\circ u'$ lies in $\JJ^2(U',Y)$. Setting
$h=-a|_{U'}\circ\operatorname{pr}_{U'}\in \JJ^2(U,Y)$, we find that $a+h$
vanishes on $U'$. Since $a+h$ is a monomorphism by hypothesis, we must have
$U'=0$, so $u$ is a split monomorphism. Replacing $p$ and $i$ by
$p\circ\alpha^{-1}$ and $\alpha\circ i$ for a suitable automorphism $\alpha$
of $E_Y$ preserves \eqref{eq:right-radical-surjection} and the factorization
property of $i$, and allows us to write
\[
 E_Y=V\oplus U,\qquad p=(f,a),\qquad i=\binom b c.
\]

\emph{Step 3: The complementary monomorphism.}
We claim that the first component $b:T\to V$ of $i$ is also a radical
monomorphism satisfying the $\JJ^2$-perturbation hypothesis. First, $b$ and $c$
are radical because $i$ is radical. If $x\in\ker b$, then $i(x)=(0,c(x))$
and $0=p(i(x))=a(c(x))$; since $a$ is injective, $c(x)=0$ and hence
$i(x)=0$, so $x=0$. Next, given any $h\in \JJ^2(T,V)$, Step~1 lets us write
$h=q\circ i$ for a radical map $q=(q_V,q_U):V\oplus U\to V$. The upper
triangular matrix
\[
 L=\begin{pmatrix}1+q_V&q_U\\0&1\end{pmatrix}
\]
is an automorphism of $V\oplus U$, and the kernel inclusion of
$p\circ L^{-1}$ is $L\circ i=\binom{b+h}{c}$. On the other hand, the
restriction of $p\circ L^{-1}$ to $U$ is
$a-f\circ(1+q_V)^{-1}\circ q_U$, which differs from $a$ by a morphism in
$\JJ^2(U,Y)$ and is therefore a monomorphism by hypothesis. As just shown for
$p$ and $b$, the injectivity of the restriction of $p\circ L^{-1}$ to $U$
implies that the first component $b+h$ of its kernel inclusion is a
monomorphism.

\emph{Step 4: Induction on the radical degree.}
We now prove by induction on $n\geq1$, simultaneously for all radical
monomorphisms $a$ satisfying the perturbation hypothesis, that
$a\circ z\in \JJ^n(W,Y)$ if and only if $z\in \JJ^{n-1}(W,U)$. Since $a$ is
radical, the ``if'' direction is immediate, as is the base case $n=1$. For
 the induction step, suppose $a\circ z\in \JJ^{n+1}(W,Y)$ with $n\geq1$. By
\eqref{eq:right-radical-surjection}, there exist $v\in \JJ^n(W,V)$ and
$u\in \JJ^n(W,U)$ such that $a\circ z=p\circ\binom{v}{u}=f\circ v+a\circ u$.
Thus $\binom{v}{u-z}$ lies in the kernel of $p$, so
$\binom{v}{u-z}=\binom{b}{c}\circ s$ for some $s:W\to T$. In particular,
$b\circ s=v\in \JJ^n(W,V)$. Applying the induction hypothesis to the
complementary monomorphism $b$ yields $s\in \JJ^{n-1}(W,T)$. Since $c$ is
radical, it follows that $u-z=c\circ s\in \JJ^n(W,U)$, and hence
$z\in \JJ^n(W,U)$, completing the induction.
\end{proof}

Applying Lemma~\ref{lem:strict-radical-layers} and its dual to the maps in
each almost split sequence yields the graded realization theorem below
(compare Iyama \cite[Theorems~5.2 and~9.2]{IyamaI} in the setting of
$\tau$-categories).

\begin{proposition}
\label{prop:graded-realization}
Let $\C=\gr_\JJ\ind(\modu A)$. Then the following statements hold:
\begin{enumerate}[label={\upshape(\arabic*)}]
\item Let $0\to\tau Y\to E_Y\to Y\to0$ be an almost split sequence in
$\modu A$, and let $\tau Y\xrightarrow{\bar\imath}E_Y\xrightarrow{\bar p}Y$
be the sequence of initial forms of its morphisms in $\add\C$. We call
$\bar p$ the \emph{right mesh map} at $Y$. For every object $W$ of $\C$, the
induced sequences
\[
 0\to\C(W,\tau Y)\to\C(W,E_Y)\to\C(W,Y),\qquad
 0\to\C(Y,W)\to\C(E_Y,W)\to\C(\tau Y,W)
\]
are exact; in particular, $\bar\imath$ is a monomorphism and $\bar p$ is an
epimorphism in $\add\C$.
\item For each indecomposable projective module $P$, the initial forms of
the components of the inclusion $\rad P\hookrightarrow P$ define a
monomorphism in $\add\C$. Dually, for each indecomposable injective module
$I$, the initial forms of the components of $I\twoheadrightarrow I/\soc I$
define an epimorphism in $\add\C$.
\item The chosen arrow representatives in $\C_1$ induce an isomorphism of
graded categories $k(\Gamma_A)\xrightarrow{\sim}\C$, where $k(\Gamma_A)$ has
the scalar mesh relations \eqref{eq:scalar-mesh}.
\end{enumerate}
\end{proposition}
\begin{proof}
Let $i:\tau Y\to E_Y$ and $p:E_Y\to Y$ be the maps in the almost split
sequence ending at $Y$. By Lemma~\ref{lem:strict-radical-layers}, $i$ induces
an injection $\C_{n-1}(W,\tau Y)\hookrightarrow\C_n(W,E_Y)$ for every
$n\geq1$. Moreover, if $w\in \JJ^n(W,E_Y)$ ($n\geq1$) satisfies
$p\circ w\in \JJ^{n+2}(W,Y)$, then \eqref{eq:right-radical-surjection} yields
some $w'\in \JJ^{n+1}(W,E_Y)$ with $p\circ w'=p\circ w$. Since
$p\circ(w-w')=0$, we have $w-w'=i\circ z$ for some $z:W\to\tau Y$, and
$i\circ z\in \JJ^n(W,E_Y)$ implies $z\in \JJ^{n-1}(W,\tau Y)$ by
Lemma~\ref{lem:strict-radical-layers}. In degree $n=0$, any
$w\in\Hom_A(W,E_Y)$ with $p\circ w\in \JJ^2(W,Y)$ differs by an element of
$\JJ(W,E_Y)$ from a morphism killed by $p$, which factors through the radical
map $i$ and therefore also lies in $\JJ(W,E_Y)$; thus $\bar p$ is injective on
$\C_0(W,E_Y)$. Since $p\circ i=0$, this proves exactness of the first
sequence in (1).

For an indecomposable projective module $P$, let $p:\rad P\hookrightarrow P$
be the inclusion. If $z\in \JJ^n(W,\rad P)$ ($n\geq0$) satisfies
$p\circ z\in \JJ^{n+2}(W,P)$, then \eqref{eq:right-radical-surjection} gives
$z'\in \JJ^{n+1}(W,\rad P)$ with $p\circ z'=p\circ z$, and injectivity of $p$
forces $z=z'\in \JJ^{n+1}(W,\rad P)$. Thus the initial form of $p$ is a
monomorphism in $\add\C$. Applying the dual arguments in the opposite module
category establishes exactness of the second sequence in (1) and the
statement for injective modules in (2).

Finally, sending each arrow of $\Gamma_A$ to its chosen class in $\C_1$
defines a graded functor $k\Gamma_A\to\C$. By \eqref{eq:scalar-mesh}, the
degree-two initial form of the relation $p\circ i=0$ is precisely the mesh
relation $r_Y$, so this functor descends to $k(\Gamma_A)\to\C$. It is full
because each radical power $\JJ^n$ is generated by products of $n$ morphisms
in $\JJ$. To see that it is faithful, we show by induction on $n\geq0$ that any
homogeneous combination $F$ of paths of length $n$ from $W$ to $Y$ whose
image vanishes in $\C_n(W,Y)$ belongs to the mesh ideal; the cases $n=0,1$
hold by definition of $\Gamma_A$. For $n\geq2$, group $F$ by its final arrow
as $F=\sum_a F_a a$, where $a:X_a\to Y$. If $Y$ is nonprojective, the
components of $\bar p:E_Y\to Y$ are $v_a a$, so the tuple
$(F_a v_a^{-1})_a$ maps to zero under $\bar p:\C_{n-1}(W,E_Y)\to\C_n(W,Y)$.
By the exactness proved in (1), there is a path combination $H$ of length
$n-2$ from $W$ to $\tau Y$ whose image under $\bar\imath$ equals the image of
$(F_a v_a^{-1})_a$ in $\C_{n-1}(W,E_Y)$. Since the $a$-component of
$\bar\imath$ is $u_a(\sigma a)$ and $u_a v_a=c_{Y,a}$, each difference
$F_a - H c_{Y,a}(\sigma a)$ has length $n-1$ and maps to zero in $\C$. By
the induction hypothesis, each such difference lies in the mesh ideal, and
hence so does
\[
 F = H r_Y + \sum_a \bigl(F_a - H c_{Y,a}(\sigma a)\bigr)a.
\]
If $Y$ is projective, then $\bar p$ is a monomorphism by (2), so each $F_a$
maps to zero in $\C_{n-1}(W,X_a)$ and lies in the mesh ideal by induction.
\end{proof}

We also record the cokernels of the right mesh maps in $\C$, which provide
minimal projective resolutions for the simple functors on $\C$.

\begin{lemma}
\label{lem:mesh-simple-tops}
For every nonprojective object $X$ of $\C$, the image of the right mesh map
$\bar p:\C(-,E_X)\to\C(-,X)$ is the radical $\bigoplus_{n\geq1}\C_n(-,X)$ of
the representable functor $\C(-,X)$, and its cokernel is the simple functor
$\C_0(-,X)$. The same conclusion holds at each projective object $P$ for the
map $\C(-,\rad P)\to\C(-,P)$ of Proposition~\ref{prop:graded-realization}(2).
\end{lemma}
\begin{proof}
Let $p:E_X\to X$ be the minimal right almost split map in $\modu A$ (or the
inclusion $\rad X\hookrightarrow X$ when $X$ is projective). By
\eqref{eq:right-radical-surjection}, $p\circ \JJ^n(W,E_X)=\JJ^{n+1}(W,X)$ for
every $n\geq0$ and every object $W$. Passing to the associated graded spaces
shows that the image of $\bar p$ on $\C(W,E_X)$ is precisely the
positive-degree ideal $\bigoplus_{n\geq1}\C_n(W,X)$. Since this ideal is
nilpotent and the degree-zero quotient has $\C_0(X,X)=k$ and $\C_0(W,X)=0$
for $W\ne X$, the positive-degree part is the radical of $\C(-,X)$ and its
cokernel is the simple functor $\C_0(-,X)$.
\end{proof}

\subsection{Evaluation on projective objects}

Following Bretscher and Gabriel \cite[Section~3.1]{BretscherGabriel}, we now
define the standard form $A_0$ as the opposite endomorphism algebra of the
projective objects in $\C$, and we show using the graded exactness properties
above that evaluating objects of $\add\C$ on the projectives gives an
equivalence $\add\C\simeq\modu A_0$.

\begin{construction}
\label{con:projective-algebra}
Let $\D$ be the full subcategory of $\C$ on the objects $P_x=e_x A$
corresponding to the indecomposable projective $A$-modules. We define the
graded algebra $A_0=\bigoplus_{x,y}A_0(x,y)$ by
\[
 A_0(x,y)=\D(P_y,P_x)
       =\bigoplus_{n\geq0}\JJ^n(P_y,P_x)/\JJ^{n+1}(P_y,P_x),
\]
with multiplication $bc=b\circ c\in\D(P_z,P_x)$ for $b\in\D(P_y,P_x)$ and
$c\in\D(P_z,P_y)$ (so that $c$ is applied first, then $b$). Equivalently,
under our path-order convention for endomorphism algebras,
$A_0=\End_{\C}(\bigoplus_x P_x)^{\op}$. The \emph{projective evaluation
functor} sends each object $X\in\add\C$ to the right $A_0$-module
$\bigoplus_x\C(P_x,X)$, on which $b\in A_0(x,y)$ acts on $m\in\C(P_x,X)$ by
$mb=m\circ b\in\C(P_y,X)$.
\end{construction}

To prove that projective evaluation is an equivalence of categories, we use a
standard homological realization criterion from the Auslander correspondence
(compare \cite[Theorem~3.3]{IyamaRealization}).

\begin{lemma}
\label{lem:projective-evaluation-criterion}
Let $\Lambda$ be a finite-dimensional algebra over a field, let $e\in\Lambda$
be an idempotent, and set $\Lambda_e=e\Lambda e$. Suppose that
\[
 \Hom_\Lambda(T,\Lambda)=\Ext^1_\Lambda(T,\Lambda)=0
 \qquad\text{for every simple $\Lambda$-module $T$ with }Te=0.
\]
Then the functor $P\mapsto Pe$ is fully faithful on $\proj\Lambda$. If, in
addition, $\operatorname{gl.dim}\Lambda\leq2$ and every finite-dimensional
injective $\Lambda_e$-module is isomorphic to $Qe$ for some
$Q\in\proj\Lambda$, then $P\mapsto Pe$ is an equivalence of categories
$\proj\Lambda\xrightarrow{\sim}\modu\Lambda_e$, and in particular
$\Lambda_e$ is representation-finite.
\end{lemma}
\begin{proof}
By induction on composition length using the long exact sequence for $\Ext$,
the assumed vanishings hold for every finite-dimensional $\Lambda$-module $T$
satisfying $Te=0$. For any right $\Lambda$-module $M$, the canonical
multiplication map $(Me)\otimes_{\Lambda_e}e\Lambda\to M$ becomes an
isomorphism upon multiplying on the right by $e$, so its kernel and cokernel
are annihilated by $e$. For any $Q\in\proj\Lambda$, applying
$\Hom_\Lambda(-,Q)$ and using the tensor--Hom adjunction therefore yields
canonical isomorphisms
\[
 \Hom_\Lambda(M,Q)\xrightarrow{\sim}
 \Hom_\Lambda((Me)\otimes_{\Lambda_e}e\Lambda,Q)
 \xrightarrow{\sim}\Hom_{\Lambda_e}(Me,Qe)
\]
whose composite sends a homomorphism $f:M\to Q$ to its restriction
$f|_{Me}:Me\to Qe$. Taking $M\in\proj\Lambda$ shows that $P\mapsto Pe$ is
fully faithful on $\proj\Lambda$.

Now suppose $\operatorname{gl.dim}\Lambda\leq2$ and every injective
$\Lambda_e$-module lies in the essential image of $\proj\Lambda$. Given any
$X\in\modu\Lambda_e$, choose the beginning of an injective coresolution
$0\to X\to I^0\to I^1$ in $\modu\Lambda_e$. By full faithfulness and the
assumption on injectives, we may lift $I^0\to I^1$ to a homomorphism
$f:Q^0\to Q^1$ in $\proj\Lambda$. Because $\operatorname{gl.dim}\Lambda\leq2$,
the kernel $\ker f$ is a second syzygy over $\Lambda$ and is therefore
projective. Since the functor $M\mapsto Me$ is exact, we obtain
$X\cong(\ker f)e$ with $\ker f\in\proj\Lambda$, proving essential
surjectivity. Finally, since $\proj\Lambda$ has only finitely many
indecomposable objects up to isomorphism, the equivalence implies that
$\Lambda_e$ is representation-finite.
\end{proof}

We now verify the hypotheses of Lemma~\ref{lem:projective-evaluation-criterion}
for $\add\C$ and $A_0$.

\begin{proposition}
\label{prop:projective-evaluation}
Let $A$ be a basic representation-finite $k$-algebra, and set
$\C=\gr_\JJ\ind(\modu A)$ and $A_0$ as in
Construction~\ref{con:projective-algebra}.
\begin{enumerate}[label={\upshape(\arabic*)}]
\item The projective evaluation functor induces an equivalence of categories
$\add\C\xrightarrow{\sim}\modu A_0$.
\item The algebra $A_0$ is basic and representation-finite.
\item If $X_0\in\modu A_0$ denotes the evaluation of an object $X\in\C$, then
\[
                         \dim_k X_0e_x=\dim_k Xe_x
\]
for every primitive idempotent $e_x$.
\end{enumerate}
\end{proposition}
\begin{proof}
Because the radical filtration on $\Hom_A(X,Y)$ is finite, passing to the
associated graded vector space preserves total dimension:
\begin{equation}
 \dim_k\C(X,Y)=\sum_{n\geq0}\dim_k \JJ^n(X,Y)/\JJ^{n+1}(X,Y)
             =\dim_k\Hom_A(X,Y).
 \label{eq:graded-hom-dimensions}
\end{equation}
In particular, $\dim_k X_0 e_x=\dim_k\C(P_x,X)=\dim_k\Hom_A(P_x,X)=\dim_k Xe_x$,
which proves (3).

To prove (1) and (2), let $G=\bigoplus_{X\in\operatorname{ob}\C}X$ and set
$\Lambda=\End_{\add\C}(G)^{\op}$. The Yoneda functor $X\mapsto R_X=\C(G,X)$
is an equivalence $\add\C\xrightarrow{\sim}\proj\Lambda$, and as $X$ ranges
over $\operatorname{ob}\C$, the modules $R_X$ form a complete set of
pairwise nonisomorphic indecomposable projective $\Lambda$-modules with
simple tops $L_X=\C_0(G,X)$. Let $e\in\Lambda$ be the sum of the idempotents
corresponding to the projective objects $P_x$. Then $e\Lambda e=A_0$, and
under the Yoneda equivalence, projective evaluation on $\add\C$ coincides
with the functor $\operatorname{ev}(M)=Me$ on $\proj\Lambda$.

\emph{Global dimension and full faithfulness.}
For each nonprojective object $X\in\C$, Proposition~\ref{prop:graded-realization}(1)
and Lemma~\ref{lem:mesh-simple-tops} yield a projective resolution in
$\modu\Lambda$:
\[
 0\longrightarrow R_{\tau X}\longrightarrow
 \bigoplus_{Y\to X}R_Y\longrightarrow R_X\longrightarrow L_X
 \longrightarrow0.
\]
For each projective object $P_x$, the same results yield a projective
resolution $0\to R_{\rad P_x}\to R_{P_x}\to L_{P_x}\to0$ of length at most
one. Consequently, every simple $\Lambda$-module has projective dimension at
most two, so $\operatorname{gl.dim}\Lambda\leq2$. Moreover, when $X$ is
nonprojective, applying $\Hom_\Lambda(-,\Lambda)$ to the resolution of $L_X$
gives by Yoneda's lemma the sequence
$0\to\C(X,G)\to\C(E_X,G)\to\C(\tau X,G)$, which is exact by
Proposition~\ref{prop:graded-realization}(1). Thus
$\Hom_\Lambda(L_X,\Lambda)=\Ext^1_\Lambda(L_X,\Lambda)=0$ for every simple
module $L_X$ annihilated by $e$, and Lemma~\ref{lem:projective-evaluation-criterion}
shows that $\operatorname{ev}$ is fully faithful on $\proj\Lambda$.

\emph{Essential surjectivity.}
Let $S_x=P_x/\rad P_x$ and $I_x=D_k(Ae_x)$ denote the simple and
indecomposable injective $A$-modules at vertex $x$. By
\eqref{eq:graded-hom-dimensions}, the evaluation $S_x^0=\operatorname{ev}(R_{S_x})$
satisfies $\dim_k S_x^0 e_y=\dim_k\Hom_A(P_y,S_x)=\delta_{xy}$, so $S_x^0$ is
the simple $A_0$-module at $x$. Set $N_x=\operatorname{ev}(R_{I_x})$. Since
$\operatorname{ev}$ is fully faithful on $\proj\Lambda$, we have
\[
 \dim_k\Hom_{A_0}(S_y^0,N_x)=\dim_k\C(S_y,I_x)=\dim_k\Hom_A(S_y,I_x)=\delta_{xy},
\]
which shows that $\soc N_x\cong S_x^0$. On the other hand,
\eqref{eq:graded-hom-dimensions} gives
\[
 \begin{aligned}
 \dim_k N_xe_y&=\dim_k\C(P_y,I_x)=\dim_k I_xe_y=\dim_k e_yAe_x\\
 &=\dim_k\C(P_x,P_y)=\dim_k e_yA_0e_x=\dim_k D_k(A_0 e_x)e_y.
 \end{aligned}
\]
Since $N_x$ has simple socle $S_x^0$, it embeds into the injective hull
$D_k(A_0 e_x)$, and equality of dimensions implies $N_x\cong D_k(A_0 e_x)$.
Thus every indecomposable injective $A_0$-module lies in the image of
$\operatorname{ev}$. Lemma~\ref{lem:projective-evaluation-criterion} now
implies that $\add\C\xrightarrow{\sim}\modu A_0$ is an equivalence. Since
the objects $P_x$ are pairwise nonisomorphic in $\C$, the algebra $A_0$ is
basic as well as representation-finite.
\end{proof}

\begin{corollary}
\label{cor:same-ar-quiver}
The Auslander--Reiten quivers of $A_0$ and $A$ are isomorphic as translation
quivers: $\Gamma_{A_0}\cong\Gamma_A$.
\end{corollary}
\begin{proof}
By Proposition~\ref{prop:projective-evaluation}(1), we may compute
$\Gamma_{A_0}$ directly inside $\add\C$. Since $\C$ is generated in degree
one, the radical of $\add\C$ is its positive-degree ideal and the radical
square is the ideal of degree at least two; hence the arrows of $\Gamma_A$
form bases of the spaces of irreducible morphisms in $\add\C$. By
Lemma~\ref{lem:mesh-simple-tops}, each right mesh map $\bar p:E_X\to X$ is
right almost split in $\add\C$. Because the indecomposable summands of $E_X$
are pairwise nonisomorphic by Lemma~\ref{lem:square-free}, any endomorphism
of $E_X$ is diagonal modulo the radical, and if it commutes with $\bar p$,
the linear independence of the arrow components modulo the radical square
forces its diagonal entries to be $1$. Thus every such endomorphism is an
automorphism, which means that $\bar p$ is right minimal. The mesh sequences
in $\add\C$ are therefore the almost split sequences of $\modu A_0$, so the
translation of $\Gamma_{A_0}$ agrees with that of $\Gamma_A$.
\end{proof}

We conclude this section with a simple example illustrating how the
categorical filtration $\JJ^n(P_y,P_x)$ used to define $A_0$ can shift the
degrees of generators of $\rad A$.

\begin{example}
\label{ex:dual-number-filtration}
Let $A=k[t]/(t^2)$ be the algebra of dual numbers. Up to isomorphism, its
indecomposable modules are $P=A$ and $S=A/(t)$, connected by the canonical
projection $\pi:P\twoheadrightarrow S$ and the inclusion $\iota:S\hookrightarrow P$
sending $1+(t)$ to $t$. In path order, the multiplication-by-$t$
endomorphism of $P$ factors as
\[
 P\xrightarrow{\pi}S\xrightarrow{\iota}P,
 \qquad \pi\iota=t,\qquad \iota\pi=0.
\]
Consequently, $\JJ$ is spanned by $\pi$, $\iota$, and $t$, while $\JJ^2$ is
spanned by $t$ and $\JJ^3=0$. On the projective endomorphism ring $A=\End_A(P)^{\op}$,
the categorical radical filtration is therefore
\[
 A\supset \JJ(P,P)=kt=\JJ^2(P,P)\supset \JJ^3(P,P)=0,
\]
whereas $(\rad A)^2=0$. The associated graded algebra $A_0$ is again
isomorphic to $k[t]/(t^2)$, but its radical generator $t$ sits in
categorical degree two, reflecting its factorization through the
nonprojective simple module $S$.
\end{example}
